\documentclass[a4paper,12pt,numbers=noenddot]{scrreprt} %Definition der KOMA-Textklasse
\input{Einstellungen}
\input{Eingaben_Benutzer}


\begin{document}
%====nicht Anfassen=======%
\automark{section} % markiert die Sections für die automatisierte Kopfzeile
\automark{chapter} % markiert die Kapitel für die automatisierte Kopfzeile
%====nicht Anfassen=======%

% Definiert ein Kapitel
\chapter{Katzen} 

%Definiert ein Unterkapitel
\section{dicke Katzen}

Oft wird angenommen, dass Katzen nicht dick seien können, sondern nur flauschig. In Abbildung \ref{fig:dicke Hauskatze} ist der Beweis dafür geliefert, dass Katzen adipös sein können.

% ==== Beispiel, wie man ein Bild einfügt ==== %
\begin{figure}[H] % das [H] muss manuell gesetzt werden, um die Position zu erzwingen
    \centering
    \includegraphics[width=0.75\linewidth]{Ressourcen/Katze-zu-dick-Sofa.jpeg}
    \caption{dicke Hauskatze}
    \label{fig:dicke Hauskatze}
\end{figure}

Allerdings ist zuviel Fleisch auf den Rippen für die meisten Hauskatzen kein Problem.

%Erzeugt ein Unterkapitel
\subsection{Hauskatzen}
Viele Hauskatzen sind meist nicht mit dem richtigen Futter ausgestattet. Viele Billigmarken, welche in Supermärkten zu erwerben sind, enthalten kaum relevante \textbf{Nährstoffe} und dienen als sogenannte \textit{Filler}, um die Katzen satt zu halten, allerdings nicht um ihnen die nötigen Nährstoffe zuzuführen.

%Definiert eine neue Seite
\newpage

\chapter{Grundlegende Befehle}
\section{Formeln}
Oft müssen auch Formel dargestellt werden. Auch hier hat Latex verschiedene Möglichkeiten. Für technische Dokumente bietet sich die Ansmath Umgebung an. Schreiben wir nun eine Formel eine bekannte Formel eines Griechen auf, so schaut diese wie folgt aus.
\begin{equation}
    a^2 + b^2 = c^2
\end{equation}
Leider wird dieser Code schnell unübersichtlich. Hier ein Beispiel für die Formel einer Hesse-Matrix einer Funktion mit zwei Variablen f(x,y).
\begin{equation}
    H =\begin{bmatrix}\frac{\partial^2 f}{\partial x^2} & \frac{\partial^2 f}{\partial x \partial y} \\\frac{\partial^2 f}{\partial y \partial x} & \frac{\partial^2 f}{\partial y^2}
    \end{bmatrix}
\end{equation}
Allerdings ist Chatgpt ein wunderbarer Freund und Helfer und kann Abhilfe schaffen für verschiedene Probleme. \cite{böswirth2012technische}

\section{Tabellen}
Auch Tabellen können sehr "leicht" eingefügt werden. Aufgrund der extremen Vielfalt ist der code für eine Tabelle eher unübersichtlich.

\begin{table}[H]
    \centering % Zentiert die Abbildung
    \begin{tabular}{|c|c|} %Definiert zwei Spalten: c = center; | = definiert den Rahmen
        \hline %definiert einen Trennstrich hier der oberste horizontale Strich der Tabelle
        \textbf{Name} & \textbf{Studiengang} \\
        \hline
        Max Mustermann & Informatik \\ %Einträge
        Anna Schmidt & Mathematik \\
        Lisa Müller & Physik \\
        Jonas Wagner & Elektrotechnik \\
        Julia Fischer & Maschinenbau \\
        \hline
    \end{tabular}
    \caption{Liste von Studenten und ihren Studiengängen}
    \label{tab:studenten_studiengaenge}
\end{table}
Hoffentlich sind diese Beispiele ausreichend für einen leichten Einstieg in die Dokumentation. Bei weiteren Fragen wendet euch bitte an mich. 
\end{document}
